% !TeX root = ../main.tex
\section{Classification of compatible \texorpdfstring{$X$}{X} measurements}
\label{appendix:x-measurements}

This appendix proves the classification of compatible $X$ measurements and derives the corresponding outcome probabilities and post-measurement states.
The proof separates measurements that enlarge the affine support from those acting within it, whose compatibility is determined by the discrete derivative of the phase polynomial.

\begingroup
\renewcommand{\label}[1]{}
\begin{xmeasurementrestatement}
    \xMeasurementClassificationStatement
\end{xmeasurementrestatement}
\endgroup
\begin{proof}
Let $\varepsilon := (-1)^\mu$ and consider the projector $\Pi_\varepsilon = \frac{1}{2}(1 + \varepsilon X_a)$.

\begin{enumerate}

\item
First, suppose that $e_a \not\in V$.
Then, $S$ and $S \oplus e_a$ are disjoint, such that
\begin{align}
    \Pi_\varepsilon \ket{\psi}
    &= \frac{1}{2\sqrt{|S|}} \left( \sum_{x\in S} \omega_8^{q(x)}\ket{x} + \varepsilon\sum_{x\in S} \omega_8^{q(x)}\ket{x \oplus e_a} \right) \\
    &= \frac{1}{2\sqrt{|S|}} \left( \sum_{x\in S} \omega_8^{q(x)}\ket{x} + \sum_{x\in S} \omega_8^{q(x) + 4\mu}\ket{x \oplus e_a} \right) \\
    &= \frac{1}{2\sqrt{|S|}} \sum_{\substack{x\in S\\ k \in\mathbb F_2}} \omega_8^{q(x) + 4k\mu} \ket{x + ke_a},
\end{align}
which remains 3PP.
The post-measurement state is given by $S'_\mu = S \cup (S \oplus e_a)$ and the unique phase polynomial $q'_\mu$ with $q'_\mu(x) = q(x)$ and $q'_\mu(x \oplus e_a) = q(x) + 4\mu$ for all $x \in S$.
Furthermore, $Pr(\mu) = |\Pi_\varepsilon\ket{\psi}|^2 = \frac{|S'_\mu|}{4|S|}|\ket{\psi}|^2 = \frac{1}{2}$.

\item
Now, suppose $e_a \in V$.
In that case,
\begin{align}
    \Pi_\varepsilon \ket{\psi}
    &= \frac{1}{2\sqrt{|S|}} \left( \sum_{x\in S} \omega_8^{q(x)}\ket{x} + \varepsilon\sum_{x\in S} \omega_8^{q(x)}\ket{x \oplus e_a} \right) \\
    &= \frac{1}{2\sqrt{|S|}} \left( \sum_{x\in S} \omega_8^{q(x)}\ket{x} + \varepsilon\sum_{x\in S} \omega_8^{q(x \oplus e_a)}\ket{x} \right) \\
    &= \frac{1}{2\sqrt{|S|}} \sum_{x\in S} \omega_8^{q(x)} (1 + \varepsilon\omega_8^{\Delta_{e_a}q(x)}) \ket{x}  \label{eqn:projected-3pp}
\end{align}
An important property of 3PP states is that all basis states in their support have equal probability.
In other words, a necessary condition for $\Pi_\varepsilon \ket{\psi}$ to remain 3PP is that  $|\bra{x}\Pi_\varepsilon \ket{\psi}| = |\bra{y}\Pi_\varepsilon \ket{\psi}|$ for all $x,y$ with $\bra{x}\Pi_\varepsilon \ket{\psi} \neq 0$ and $\bra{y}\Pi_\varepsilon \ket{\psi} \neq 0$.
Following \cref{eqn:projected-3pp}, this is equivalent to
\[
    \left|1 + \varepsilon\omega_8^{\Delta_{e_a}q(x)}\right| = \left|1 + \varepsilon\omega_8^{\Delta_{e_a}q(y)}\right|
\]
for all $x,y$ with $1 + \varepsilon\omega_8^{\Delta_{e_a}q(x)} \neq 0$ and $1 + \varepsilon\omega_8^{\Delta_{e_a}q(y)} \neq 0$.
Computing those magnitudes, we get:
\begin{align*}
    \left| 1 + \varepsilon\omega_8^{\Delta_{e_a}q(x)} \right| &=
    \left| \omega_8^{\Delta_{e_a} q(x) / 2} \left( \omega_8^{-\Delta_{e_a} q(x) / 2} + \varepsilon\omega_8^{\Delta_{e_a} q(x) / 2} \right) \right| \\
    &= \begin{cases}
        2 \left| \cos \left(\dfrac{\Delta_{e_a} q(x) \pi}{8} \right) \right| & \text{if } \varepsilon = +1, \\
        2 \left| \sin \left(\dfrac{\Delta_{e_a} q(x) \pi}{8} \right) \right| & \text{if } \varepsilon = -1.
    \end{cases}
\end{align*}
For $\varepsilon = +1$, this means that the image of $\Delta_{e_a}q$ must be a subset of one of $\{0, 4\}, \{2,4,6\}, \{1,4,7\}$, or $\{3,4,5\}$.
Similarly, for $\varepsilon = -1$ its image must be a subset of one of $\{0, 4\}, \{0,2,6\}, \{0,1,7\}$, or $\{0,3,5\}$.
Since both conditions must hold simultaneously, compatibility requires
\begin{align} \label{eqn:gradient-two-valued}
    \operatorname{Im}(\Delta_{e_a}q)\subseteq D
    \quad\text{for some}\quad
    D\in\bigl\{\{0,4\},\{2,6\},\{1,7\},\{3,5\}\bigr\}.
\end{align}
We will discuss these cases separately:

\begin{enumerate}

\item
Suppose $\Delta_{e_a} q(x) \in \{ 0, 4 \}$ for all $x \in S$, i.e. $\Delta_{e_a} q(x) = 4f(x)$ for some function $f : S \to \mathbb F_2$.
Then,
\[
    1 + \varepsilon\omega_8^{E(x)}
    = \begin{cases}
        2 & \text{if } f(x) = \mu \\
        0 & \text{otherwise}.
    \end{cases}
\]
Inserting this into \cref{eqn:projected-3pp}, we see that basis states with $f(x) \neq \mu$ are removed from the support while the others are scaled by 2, giving us
\[ \Pi_\varepsilon\ket{\psi} = \frac{1}{\sqrt{|S|}} \sum_{\substack{x \in S \\ f(x) = \mu}} \omega_8^{q(x)} \ket{x}. \]
This state remains 3PP for both outcomes $\mu = 0,1$ simultaneously iff $f$ is affine.
In this case, we can read off that $S'_\mu = S \cap \{ x ~|~ f(x) = \mu \}$ and $q'_\mu = q$.
Furthermore, $Pr(\mu) = |\Pi_\varepsilon\ket{\psi}|^2 = \frac{|S'_\mu|}{|S|}|\ket{\psi}|^2 = \frac{|S'_\mu|}{|S|}$.
Since $S$ and $S'_\mu$ are affine spaces differing by a single constraint, this ratio must be either 0, 1, or $\frac{1}{2}$ depending on whether $f(x) = \mu$ holds for none, all, or some $x \in S$.


\item
Suppose $\Delta_{e_a}q(x) \in \{2, 6\}$ for all $x \in S$, i.e. $\Delta_{e_a}q(x) = 2 + 4g(x) \pmod 8$ for some function $g: S \to \mathbb F_2$.
Then,
\[
    1 + \varepsilon\omega_8^{\Delta_{e_a}q(x)}
    = 1 + i\varepsilon(-1)^{g(x)}
    = \sqrt 2 \omega_8^{\varepsilon(-1)^{g(x)}}
    = \sqrt 2 \omega_8^{\varepsilon(1 - 2g(x))},
\]
where we used the identities $1+s i = \sqrt{2} e^{is\pi/4}$ for $s \in \{-1,1\}$ and  $(-1)^m = 1 - 2m$ for $m \in \{0,1\}$.
Inserting this into \cref{eqn:projected-3pp}, we get the projected state
\[ \Pi_\varepsilon \ket{\psi} = \frac{\sqrt 2 \omega_8^\varepsilon}{2\sqrt{|S|}} \sum_{x \in S} \omega_8^{q(x) - 2\varepsilon g(x)} \ket{x} \]
which remains 3PP iff $2g$ is a valid third-order phase polynomial.
In this case, we can read off that $S' = S$ and $q'(x) = q(x) - 2\varepsilon g(x)$.
Furthermore, $Pr(\mu) = |\Pi_\varepsilon\ket{\psi}|^2 = \frac{1}{2}|\ket{\psi}|^2 = \frac{1}{2}$.


\item
Suppose $\Delta_{e_a}q(x) \in \{ d, -d \}$ for all $x \in S$ where $d \in \{ 1, 3, 5, 7 \}$, i.e. $\Delta_{e_a}q(x) = (-1)^{h(x)}d \pmod 8$ for some function $h: S \to \mathbb F_2$.
Then,
\begin{align*}
    1 + \varepsilon\omega_8^{\Delta_{e_a}q(x)}
    &= (1 + \varepsilon\omega_8^d) \left(\frac{1 + \varepsilon\omega_8^{(-1)^{h(x)} d}}{1 + \varepsilon\omega_8^{(-1)^{h(x) + 1} d}}\right)^{h(x)} \\
    &= (1 + \varepsilon\omega_8^d) (\varepsilon\omega_8^{-d})^{h(x)} \\
    &= (1 + \varepsilon\omega_8^d) \omega_8^{(4\mu - d)h(x)}.
\end{align*}
Inserting this into \cref{eqn:projected-3pp}, we get the projected state
\[ \Pi_\varepsilon \ket{\psi} = \frac{1 + \varepsilon\omega_8^d}{2\sqrt{|S|}} \sum_{x \in S} \omega_8^{q(x) + (4\mu - d)h(x)} \ket{x}. \]
We claim that this state remains 3PP iff $h$ is affine over $\mathbb F_2$.
First, suppose $h$ is affine over $\mathbb F_2$, i.e. $h(x) = a_0 \oplus \left(\bigoplus_i a_i x_i\right)$ for some $a_0,a_i \in \mathbb F_2$. Writing $a=(a_i)_i$, we have
\[
    h(x) = a_0 \oplus \langle a,x\rangle
    = a_0 + (1-2a_0)\langle a,x\rangle \pmod 8.
\]
Thus, $h$ has the form of \cref{eq:phase-polynomial-sum-parities}, and hence $q'_\mu=q+(4\mu-d)h$ is a third-order phase polynomial.
Conversely, note that
\[ h(x) = (4\mu - d)^{-1}(q'_\mu(x) - q(x)) \pmod 2. \]
Thus, if $q'_\mu$ is a 3PP phase polynomial, then $(4\mu - d)^{-1}(q'_\mu - q)$ must also be one.
Reducing mod 2 removes all quadratic and cubic terms, showing that $h$ is affine over $\mathbb F_2$.
If the measurement is compatible, we can furthermore read off that $S'_\mu = S$ and
\[ Pr(\mu) = |\Pi_\varepsilon\ket{\psi}|^2 = \left|\frac{1 + \varepsilon\omega_8^d}{2}\right|^2 |\ket{\psi}|^2 = \frac{1}{2}(1 + \varepsilon\cos(d\pi/4)). \]

\end{enumerate}

In any other case, the measurement cannot be compatible as established in \cref{eqn:gradient-two-valued}. \qedhere

\end{enumerate}
\end{proof}

The measurement classification in \cref{thm:x-measurements} depends on the derivative $\Delta_{e_a}q$ of the phase polynomial in the direction of the measured qubit.
Remember that, by \cref{prop:derivative-from-E},  we can directly construct this derivative from the DCP stabilizer group of $\ket{\psi}$ without having to explicitly track the phase polynomial itself.

More precisely, let $\ket{\psi}$ be a 3PP state with phase polynomial $q$ and DCP stabilizer group $G$. Suppose $e_a \in X(G)$ and choose an element $g \in G$ with $\pi(g) = e_a$.

Write $g=X_a\delta$, where $\delta=\omega_8^kD(s,M)$ is diagonal.
The stabilizer equation $g\ket{\psi}=\ket{\psi}$ gives
\begin{align}
    \Delta_{e_a}q(x)
    =k+2\sum_i s_i x_i+4\sum_{i<j}M_{ij}x_ix_j
    \pmod 8
\end{align}
for every $x\in S$.
Thus, we can extract $\Delta_{e_a} q$ directly from the normal form of the DCP stabilizer element $g$.
Furthermore, given $\Delta_{e_a} q$, we can distinguish the different cases of \cref{thm:x-measurements} efficiently.
To this end, Let $S=x_0\oplus V$, and let $v^{(1)},\ldots,v^{(r)}$ be a basis of $V$.
Define the restricted derivative
\begin{align}
    \widetilde D(c)
    :=\Delta_{e_a}q\left(x_0\oplus\bigoplus_{i=1}^r c_i v^{(i)}\right),
    \qquad c\in\mathbb F_2^r.
\end{align}
By affine pullback, this is a Clifford phase polynomial, which we write as
\begin{align}
    \widetilde D(c)
    =\tilde k+2\sum_i\tilde s_i c_i
    +4\sum_{i<j}\tilde M_{ij}c_ic_j
    \pmod 8,
\end{align}
where $\tilde k\in\mathbb Z_8$, $\tilde s_i\in\mathbb Z_4$, and $\tilde M$ is symmetric over $\mathbb F_2$ with zero diagonal.

\begin{proposition}\label{prop:x-measurement-case-decision}
    The different cases of \cref{thm:x-measurements} can be decided based on $\tilde k$, $\tilde s$, and $\tilde M$:
    \begin{enumerate}
        \item
        Case \ref{thm:x-measurements:even-0-4} applies iff $\tilde k \in \{0, 4\}$, $\tilde s_i \in \{0, 2\}$, and $\tilde M_{ij} = 0$.

        \item
        Case \ref{thm:x-measurements:even-2-6} applies iff $\tilde k \in \{2, 6\}$, $\tilde s_i \in \{0, 2\}$, and $\rank \tilde M \leq 2$.

        \item
        Case \ref{thm:x-measurements:odd} applies iff $\tilde k$ is odd,
        $\tilde s_i\in\{0,-\tilde k\}\pmod 4$ for every $i$, and
        $\tilde M_{ij}=\tilde s_i\tilde s_j\pmod 2$ for every $i<j$.
    \end{enumerate}
\end{proposition}
\begin{proof}
We begin by giving explicit representations for $\tilde k, \tilde s_i$, and  $\tilde M_{ij}$ in terms of $\Delta_{e_a} q$.
Concretely, noting that $\widetilde D(0) = \tilde k$, $\widetilde D(e_i) = \tilde k + 2\tilde s_i$, and $\widetilde D(e_i \oplus e_j) = \tilde k + 2\tilde s_i + 2\tilde s_j + 4\tilde M_{ij}$, we have
\begin{align}
    \tilde k &= \Delta_{e_a}q(x_0) \pmod 8 \label{eqn:k-tilde} \\
    2\tilde s_i &= \Delta_{e_a}q(x_0 \oplus v^{(i)}) - \Delta_{e_a}q(x_0) \pmod 8 \label{eqn:s-tilde} \\
    4\tilde M_{ij} &= \Delta_{e_a}q(x_0 \oplus v^{(i)} \oplus v^{(j)}) - \Delta_{e_a}q(x_0 \oplus v^{(i)}) - \Delta_{e_a}q(x_0 \oplus v^{(j)}) + \Delta_{e_a}q(x_0) \pmod 8 \label{eqn:M-tilde}.
\end{align}
In other words $2\tilde s_i = \Delta_{v^{(i)}}\Delta_{e_a}q(x_0)$ and $4\tilde M_{ij} = \Delta_{v^{(i)}}\Delta_{v^{(j)}}\Delta_{e_a}q(x_0)$ are the first and second discrete derivatives of $\Delta_{e_a}q$ at $x_0$ in the directions of $V$'s basis vectors.
Furthermore, we note the following facts about arithmetic over Boolean variables $a_i \in \mathbb F_2$ modulo 8:
\begin{align}
    2\bigoplus a_i &= 2\sum_i a_i - 4\sum_{i<j} a_ia_j \pmod 8 \label{eqn:2-xor} \\
    4\bigoplus a_i &= 4\sum_{i} a_i \pmod 8 \label{eqn:4-xor}
\end{align}
We consider each case separately:
\begin{enumerate}
\item
First, suppose that $\Delta_{e_a}\tilde q(c) = 4f(c) \pmod 8$ for some affine function $f: \mathbb F_2^r \to \mathbb F_2$. Write $f(c) = \bigoplus_i f_ic_i \oplus f_0$.
Using \cref{eqn:4-xor}, we see that $\Delta_{e_a}\tilde q(c) = 4f_0 + 4\sum_i f_ic_i \pmod 8$.
We can read off that $\tilde k := 4f_0 \in \{0, 4\}$, $2\tilde s_i := 4f_i$ is even and $\tilde M_{ij} = 0$ since there are no quadratic terms.

Conversely, suppose $\tilde k \in \{0,4\}$, $\tilde s_i \in \{0,2\}$, and $\tilde M_{ij} = 0$.
Then $f(c) := \left(\bigoplus_i \frac{s_i}{2}c_i\right) \oplus \frac{k}{4}$ is affine and $\widetilde D(c) = 4f(c)$ such that case \ref{thm:x-measurements:even-0-4} applies.

\item
Suppose $\Delta_{e_a}\tilde q(c) = 2 + 4g(c) \pmod 8$ for some function $g: \mathbb F_2^r \to \mathbb F_2$ such that $2g$ is a 3PP phase polynomial.
Evaluating equations \eqref{eqn:k-tilde} and \eqref{eqn:s-tilde}, we get
\begin{align*}
    \tilde k &= \Delta_{e_a}\tilde q(0) = 2 + 4g(0) \pmod 8 \\
    2\tilde s_i &= \Delta_{e_a}\tilde q(e_i) - \Delta_{e_a}\tilde q(0) = 4(g(e_i) - g(0)) \pmod 8
\end{align*}
showing that $\tilde k \in \{2,6\}$ and that all $\tilde s_i$ are even.
Then
\begin{align*}
    4g(c)
    = \Delta_{e_a}\tilde q(c) - 2
    = \tilde k - 2 + 2\sum_i \tilde s_ic_i + 4\sum_{i<j} \tilde M_{ij}c_ic_j \pmod 8
\end{align*}
Using \cref{eqn:4-xor}, we conclude that
\begin{equation} \label{eqn:g-quadratic}
    g(c) = \frac{\tilde k - 2}{4} \oplus \left(\bigoplus_i \frac{\tilde s_i}{2}c_i\right) \oplus \left(\bigoplus_{i<j} \tilde M_{ij}c_ic_j\right).
\end{equation}
In other words, $g$ must be a quadratic Boolean function.
Therefore, $\rank \tilde M \leq 2$ by Lemma~\ref{lem:3pp-rank}.

Conversely, suppose $\tilde k \in \{2, 6\}$, $\tilde s_i \in \{0, 2\}$, and $\rank\tilde M \leq 2$.
Define $g$ by \cref{eqn:g-quadratic} and use \cref{eqn:4-xor} to show
\begin{align}
    2+4g(c)
    =\tilde k+2\sum_i\tilde s_i c_i
    +4\sum_{i<j}\tilde M_{ij}c_ic_j
    =\widetilde D(c)
    \pmod 8.
\end{align}
By Lemma~\ref{lem:3pp-rank}, $2g$ is a 3PP phase polynomial such that case \ref{thm:x-measurements:even-2-6} applies.


\item
Suppose first that case \ref{thm:x-measurements:odd} applies.
By absorbing the constant term of the affine function into the sign of $d$, we may write
\begin{align}
    \widetilde D(c)=\tilde k(-1)^{h(c)},
    \qquad
    h(c)=\bigoplus_i h_i c_i,
\end{align}
where $\tilde k$ is odd and $h_i\in\mathbb F_2$.
Expanding the parity modulo $8$ gives
\begin{align}
    \widetilde D(c)
    &=\tilde k-2\tilde k\sum_i h_i c_i
    +4\tilde k\sum_{i<j}h_ih_jc_ic_j
    \pmod 8.
\end{align}
Comparing coefficients yields
\begin{align}
    \tilde s_i&=-\tilde k h_i\pmod 4,\\
    \tilde M_{ij}&=h_ih_j\pmod 2.
\end{align}
Since $\tilde k$ is odd, we also have $h_i=\tilde s_i\pmod 2$, proving the stated conditions.

Conversely, suppose those conditions hold.
Define $h_i:=\tilde s_i\pmod 2\in\mathbb F_2$ and $h(c):=\bigoplus_i h_ic_i$.
The condition $\tilde s_i\in\{0,-\tilde k\}\pmod 4$ implies $\tilde s_i=-\tilde k h_i\pmod 4$, while the quadratic condition gives $\tilde M_{ij}=h_ih_j\pmod 2$.
Therefore,
\begin{align}
    \tilde k(-1)^{h(c)}
    &=\tilde k-2\tilde k\sum_i h_i c_i
    +4\tilde k\sum_{i<j}h_ih_jc_ic_j\\
    &=\tilde k+2\sum_i\tilde s_i c_i
    +4\sum_{i<j}\tilde M_{ij}c_ic_j\\
    &=\widetilde D(c)\pmod 8.
\end{align}
Thus case \ref{thm:x-measurements:odd} applies.
\qedhere
\end{enumerate}
\end{proof}


\begin{lemma} \label{lem:3pp-rank}
    Let $g: \mathbb F_2^n \to \mathbb F_2$ be a quadratic Boolean function given by a bilinear form
    \[ g(x) = g_0 \oplus \left( \bigoplus_i a_i x_i \right) \oplus \left( \bigoplus_{i<j} M_{ij}x_ix_j \right) \]
    with $g_0 \in \mathbb F_2, a\in\mathbb F_2^n, M \in \mathbb F_2^{n\times n}$ where $M$ is symmetric with zero diagonal.
    Then $2g$ is a 3PP phase polynomial iff $\rank M \leq 2$.
\end{lemma}
\begin{proof}
We first put the quadratic part into canonical coordinates, then determine when the weighted polynomial for $2g$ contains quartic terms.
Since $M$ is symmetric with zero diagonal, its rank is even, say $\operatorname{rank}M=2t$.
An invertible linear change of coordinates puts $g$ into the form
\begin{align}
    \hat g(z)
    =\ell(z)\oplus\bigoplus_{i=1}^t z_{2i-1}z_{2i},
    \qquad z\in\mathbb F_2^n,
\end{align}
where $\ell$ is affine and may depend on all $n$ variables.
By affine pullback in both directions, $2g$ is a third-order phase polynomial iff $2\hat g$ is one.

Expanding the XOR modulo $8$ gives
\begin{align}
    2\hat g(z)
    &=2\ell(z)+2\sum_{i=1}^t z_{2i-1}z_{2i}
    -4\sum_{i=1}^t\ell(z)z_{2i-1}z_{2i}\\
    &\quad-4\sum_{i<j}z_{2i-1}z_{2i}z_{2j-1}z_{2j}
    \pmod 8.
\end{align}
The first term is a third-order phase polynomial because $\ell$ is affine.
The second term has quadratic coefficients divisible by $2$.
In the third term, the prefactor $4$ allows us to expand $\ell$ over $\mathbb F_2$, so this term has degree at most three and all coefficients are divisible by $4$.
Thus, when $t\leq 1$, every term has the required weighted form.

If $t\geq 2$, the last sum contains $z_1z_2z_3z_4$ with coefficient $4$ modulo $8$.
None of the other terms has degree four, so this coefficient cannot cancel.
By uniqueness of the multilinear polynomial representation modulo $8$, $2\hat g$ is then not a third-order phase polynomial.
Consequently, $2g$ is a third-order phase polynomial iff $t\leq 1$, equivalently $\operatorname{rank}M\leq 2$.
\end{proof}
